% Options for packages loaded elsewhere
\PassOptionsToPackage{unicode}{hyperref}
\PassOptionsToPackage{hyphens}{url}
\PassOptionsToPackage{dvipsnames,svgnames*,x11names*}{xcolor}
%
\documentclass[
  11pt,
]{article}
\usepackage{lmodern}
\usepackage{amssymb,amsmath}
\usepackage{ifxetex,ifluatex}
\ifnum 0\ifxetex 1\fi\ifluatex 1\fi=0 % if pdftex
  \usepackage[T1]{fontenc}
  \usepackage[utf8]{inputenc}
  \usepackage{textcomp} % provide euro and other symbols
\else % if luatex or xetex
  \usepackage{unicode-math}
  \defaultfontfeatures{Scale=MatchLowercase}
  \defaultfontfeatures[\rmfamily]{Ligatures=TeX,Scale=1}
\fi
% Use upquote if available, for straight quotes in verbatim environments
\IfFileExists{upquote.sty}{\usepackage{upquote}}{}
\IfFileExists{microtype.sty}{% use microtype if available
  \usepackage[]{microtype}
  \UseMicrotypeSet[protrusion]{basicmath} % disable protrusion for tt fonts
}{}
\makeatletter
\@ifundefined{KOMAClassName}{% if non-KOMA class
  \IfFileExists{parskip.sty}{%
    \usepackage{parskip}
  }{% else
    \setlength{\parindent}{0pt}
    \setlength{\parskip}{6pt plus 2pt minus 1pt}}
}{% if KOMA class
  \KOMAoptions{parskip=half}}
\makeatother
\usepackage{xcolor}
\IfFileExists{xurl.sty}{\usepackage{xurl}}{} % add URL line breaks if available
\IfFileExists{bookmark.sty}{\usepackage{bookmark}}{
\usepackage{hyperref}
}
\hypersetup{
  pdftitle={Web et Interaction Homme Machine, programmation côté client},
  pdfauthor={Première NSI, Lycée du Parc},
  colorlinks=true,
  linkcolor=Maroon,
  filecolor=Maroon,
  citecolor=Blue,
  urlcolor=Blue,
  pdfcreator={LaTeX via pandoc}}
\urlstyle{same} % disable monospaced font for URLs
\usepackage[top=20mm,left=20mm,right=20mm,heightrounded]{geometry}
\usepackage{listings}
\newcommand{\passthrough}[1]{#1}
\lstset{defaultdialect=[5.3]Lua}
\lstset{defaultdialect=[x86masm]Assembler}
\usepackage{longtable,booktabs}
% Correct order of tables after \paragraph or \subparagraph
\usepackage{etoolbox}
\makeatletter
\patchcmd\longtable{\par}{\if@noskipsec\mbox{}\fi\par}{}{}
\makeatother
% Allow footnotes in longtable head/foot
\IfFileExists{footnotehyper.sty}{\usepackage{footnotehyper}}{\usepackage{footnote}}
\makesavenoteenv{longtable}
\usepackage{graphicx}
\makeatletter
\def\maxwidth{\ifdim\Gin@nat@width>\linewidth\linewidth\else\Gin@nat@width\fi}
\def\maxheight{\ifdim\Gin@nat@height>\textheight\textheight\else\Gin@nat@height\fi}
\makeatother
% Scale images if necessary, so that they will not overflow the page
% margins by default, and it is still possible to overwrite the defaults
% using explicit options in \includegraphics[width, height, ...]{}
\setkeys{Gin}{width=\maxwidth,height=\maxheight,keepaspectratio}
% Set default figure placement to htbp
\makeatletter
\def\fps@figure{htbp}
\makeatother
\setlength{\emergencystretch}{3em} % prevent overfull lines
\providecommand{\tightlist}{%
  \setlength{\itemsep}{0pt}\setlength{\parskip}{0pt}}
\setcounter{secnumdepth}{5}

\title{Web et Interaction Homme Machine, programmation côté client}
\author{Première NSI, \href{https://frederic-junier.org/}{Lycée du
Parc}}
\date{}

%%%jolis boites

\usepackage{fancybox, graphicx}



%%%%%%%%%%%%%%%%Packages et Macros Frederic%%%%%%%%%%%%%%%%%%%%%%%%%%%%%


%%%%Insertion de liens hypertextes %%%%

            
%%%%%%%%%%PSTricks%%%%%%%%%%%%

\usepackage{pstricks,pst-plot,pst-text,pst-tree,pst-eps,pst-fill,pst-node,pst-math,pstricks-add,pst-xkey,pst-eucl}


%%%%%%%Tikz%%%%%%%%%%%%%%%
\usepackage{pgf,tikz,tkz-tab}
% Pour les tableaux de signes ou de variations avec tkz-tab voir https://zestedesavoir.com/tutoriels/439/des-tableaux-de-variations-et-de-signes-avec-latex/#1-13389_tikz-un-package-qui-en-a-dans-le-ventre
\usetikzlibrary{arrows}
\usetikzlibrary{shapes.geometric}
\usetikzlibrary{shapes.geometric}
\usetikzlibrary{petri}
\usetikzlibrary{decorations}
\usetikzlibrary{arrows}
\usetikzlibrary{math}
 %Variables must be declared in a tikzmath environment but
       % can be used outside
%       \tikzmath{int \n; \n = 508; \x1 = 1; \y1 =1; 
%                   %computations are also possible
%                    \x2 = \x1 + 1; \y2 =\y1 +3; } 


%%%%%%%%%%%%%%%%%%%%%%%%%%%%%%%%%%%%%%%%
%%%%%%%%%%%Commandes Tikz Perso%%%%%%%%%%%%%%%

% Définition des nouvelles options xmin, xmax, ymin, ymax
% Valeurs par défaut : -3, 3, -3, 3
\tikzset{
xmin/.store in=\xmin, xmin/.default=-3, xmin=-3,
xmax/.store in=\xmax, xmax/.default=3, xmax=3,
ymin/.store in=\ymin, ymin/.default=-3, ymin=-3,
ymax/.store in=\ymax, ymax/.default=3, ymax=3,
}
% Commande qui trace la grille entre (xmin,ymin) et (xmax,ymax)
\newcommand {\grille}[2]
{\draw[help lines,black, thick] (\xmin,\ymin) grid[xstep=#1, ystep=#2] (\xmax,\ymax);}
% Commande \axes
\newcommand {\axes} {
\draw[->,very thick] (\xmin,0) -- (\xmax,0);
\draw[->,very thick] (0,\ymin) -- (0,\ymax);
\draw (0.95*\xmax, 0) node[above] {};
\draw (0, 0.95*\ymax) node[left] {};
}
% Commande qui limite l?affichage à (xmin,ymin) et (xmax,ymax)
\newcommand {\fenetre}
{\clip (\xmin,\ymin) rectangle (\xmax,\ymax);}

%Exemple d'utilisation

%\begin{center}
%\begin{tikzpicture} [xmin=-2,xmax=2,ymin=0,ymax=5]
%\grille{1} \axes \fenetre
%\draw plot[smooth] (\x,\x^2);
%\end{tikzpicture}
%\end{center}

%style pour la perspective cavalière française
%voir Tikz pour l'impatient page 68
\tikzset{math3d/.style=
{x= {(-0.353cm,-0.353cm)}, z={(0cm,1cm)},y={(1cm,0cm)}}}

%%%%%%%Symbole pour code calculatrice%%%%%%

%Flèche remplie pour défilement de menu

\newcommand{\flechefillright}{
\begin{tikzpicture}[scale=0.15] \fill (0,0)--(2,1)--(0,2)--cycle;
\end{tikzpicture}}

%%%%%%%%%%%%%Symboles pour calculatrice Casio%%%%
\newcommand{\execasio}{\Pisymbol{psy}{191}} %Retour chariot
\newcommand{\dispcasio}{\begin{pspicture}(.1,.1)\pspolygon*(.1,0)(.1,.1)\end{pspicture}} %Triangle « Disp »
\newcommand{\dispcasiotikz}{\begin{tikzpicture}[scale=0.2]
\fill (0,0) -- (1,0) -- (1,1) -- cycle;
\end{tikzpicture}} %Triangle « Disp »
%

%Fleche entre deux lignes, d'apres 'un bon petit' : http://forum.mathematex.net/latex-f6/fleches-entre-deux-lignes-pour-resolution-d-equation-t10283.html#p99817
\newcommand\addnode[1]{\Rnode{#1}{}}
\newcommand\linknode[3]{\ncbar[angleA=0,angleB=0,nodesep=1ex,arm=10ex,offset=-2pt]{->}{#1}{#2}\Aput{\vphantom{x}#3}}


%%Commande pour touche de calculatrice

\newcommand\tc[1]{%
{
\begin{tikzpicture}
\node[draw,rectangle,rounded corners=3pt] (P) at (0,0){#1};
\end{tikzpicture}
}
}

%%%%%%%%%%%%%%%%%%%%%%%%%%%%%%%%%%%%%%%%
%%%%%%%%%%%Fin Commandes Tikz%%%%%%%%%%%%%%%


%%%%%%%%%%%%Specifiques%%%%%%%%%%%
\usepackage{wrapfig}
%pour insérer une figure à droite ou à gauche d'un texte
%\begin{wrapfigure}[nb lignes]{placement l,r,c,i(inside),o(outside)}[overhang]{width}
%ce package fonctionne mal à proximité des listes
%%%%%%%%%%%%%%%%%%%%%%%%%%%%%%%%%%%%%

%%%%%Environnements et symboles spéciaux pour faire joli%%%%%%

%%%Bclogo, pour des environnements + jolis avec insertion de logo%%%%
%Dépendances de  bclogo
\usepackage{xkeyval}  
\usepackage{etoolbox}
\usepackage{ifpdf}
\usepackage[framemethod=tikz]{mdframed}
\usepackage[tikz]{bclogo}

%\newcommand\bcpython{\includegraphics[width=17pt]{/home/fjunier/Maths/python-logo.eps}}
\newcommand\bcpython{\includegraphics[width=17pt]{/home/fjunier/Maths/python-logo.png}}
%\newcommand\bcpython{\includegraphics[width=17pt]{/home/frederic/Maths/python-logo.png}}

%% Framed
\usepackage{framed}  %Le package « framed» Crée 3 nouveaux environnements, qui se comportent comme des minipage de largeur \linewidth, mais permettant en plus de se casser entre plusieurs pages.     * framed : avec un cadre autour;     * shaded : avec un fonc coloré (il faut définir la couleur shadecolor);     * leftbar : avec une barre le long du côté gauche.

%%%%%%%%%%%%%%%%%%%Présentation de codes sources%%%%%%%%%%%%%%%%%
\usepackage{listings}
%On utilise l?environnement lstlisting pour insérer
%un code source.
%En plus de l?environnement lstlisting, on peut également utiliser la
%commande \lstinline qui fonctionne comme la commande \verb, en ce
%sens qu?on peut utiliser n?importe quel caractère comme délimiteur. Enfin,
%la commande \lstinputlisting permet de charger un code source depuis
%un fichier externe.
%Il y a deux manières de préciser des options : soit via l?option de l?envi-
%ronnement ou de la commande, soit en utilisant la commande \lstset
%qui permet de définir des options de manière globale.

\lstset{ %
  language=Python,                % the language of the code
  basicstyle=\ttfamily,           % the size of the fonts that are used for the code
  %numbers=left,                   % where to put the line-numbers
  numberstyle=\tiny,  % the style that is used for the line-numbers
  %stepnumber=2,                   % the step between two line-numbers. If it's 1, each line 
                                  % will be numbered
  %numbersep=5pt,                  % how far the line-numbers are from the code
  backgroundcolor=\color{white},      % choose the background color. You must add \usepackage{color}
  showspaces=false,               % show spaces adding particular underscores
  showstringspaces=false,         % underline spaces within strings
  showtabs=false,                 % show tabs within strings adding particular underscores
  frame=single,                   % adds a frame around the code
  rulecolor=\color{black},        % if not set, the frame-color may be changed on line-breaks within not-black text (e.g. comments (green here))
  tabsize=4,                      % sets default tabsize to 2 spaces
  captionpos=b,                   % sets the caption-position to bottom
  breaklines=true,                % sets automatic line breaking
  breakatwhitespace=false,        % sets if automatic breaks should only happen at whitespace
  %title=\lstname,                   % show the filename of files included with \lstinputlisting;
                                  % also try caption instead of title
  breakindent=1cm,
  keywordstyle=\color{blue},          % keyword style
  commentstyle=\color{red},       % comment style
  %stringstyle=\ttfamily\color{green},         % string literal style
  escapeinside={\%*}{*)},            % if you want to add LaTeX within your code
  morekeywords={*,...},              % if you want to add more keywords to the set
  deletekeywords={...}              % if you want to delete keywords from the given language
  upquote=true,columns=flexible,
xleftmargin=1cm,xrightmargin=1cm,
 inputencoding=utf8,			%Les lignes qui suivent sont pour le codage utf8
  extendedchars=true,
  literate=%
            {é}{{\'{e}}}1
            {è}{{\`{e}}}1
            {ê}{{\^{e}}}1
            {ë}{{\¨{e}}}1
            {û}{{\^{u}}}1
            {ù}{{\`{u}}}1
            {â}{{\^{a}}}1
            {à}{{\`{a} }}1
            {î}{{\^{i}}}1
            {ô}{{\^{o}}}1
            {ç}{{\c{c}}}1
            {Ç}{{\c{C}}}1
            {É}{{\'{E}}}1
            {Ê}{{\^{E}}}1
            {À}{{\`{A}}}1
            {Â}{{\^{A}}}1
            {Î}{{\^{I}}}1
}

\lstdefinestyle{rond}{
  numbers=none,
  backgroundcolor=\color{gristclair},
  frameround =tttt
}

\lstdefinestyle{compil}{
  numbers=none,
  backgroundcolor=\color{gristclair}
}
%\lstset{language=Python,basicstyle=\small , frame=single,tabsize=4,showspaces=false,showtabs=false,showstringspaces=false,numbers=left,numberstyle=\tiny , extendedchars=true}



%%%%%%%%%%%%%%%%%%%%%%%%%%%%%%%%%%%%%%%%%%%%%%%%%%%%%%%%%%%%%%%%%%%%%%%%
%%%%%%%%%%%%%%%%%%%%Environnements persos%%%%%%%%%%%%%%%%%%%%%%%%%%%%%%%%
%Syntaxe :
%\newenvironment{nom}[nombre d'args][defaut]{definitions initiales}{definitions finales}
%definitions intiales sont les commandes appelées par \begin{nom}
%Definitions finales sont les commandes appelées par \end{nom}

%%%%%%%%%%%%%%%%Définitions des environnemts persostheoreme, exemple ..%%%%
%%%% Exercice avec encadré %%%%
\newcounter{exo}
\newenvironment{exercice}[1]
{\par \medskip   \addtocounter{exo}{1} \noindent  
\begin{bclogo}[arrondi =0.1,   noborder = true, logo=\bccrayon, marge=4]{~\textbf{Exercice} \textbf{\theexo} {\itshape #1} }  \par}
{
\end{bclogo}
 \par \bigskip }

%%Axiomes, Theoremes, Propriété, Définition, Methode, Preuve


\newenvironment{axiome}[1]
{\par \medskip   \begin{leftbar} \noindent \underline{\textbf{Axiome}}\hspace{0.5cm}{\itshape #1}   \vspace*{10pt} \par }
{\end{leftbar}  \par \medskip }


\newcounter{thme}
\newenvironment{theoreme}[1]
{\par \medskip  \addtocounter{thme}{1} \noindent  
\begin{bclogo}[arrondi =0.1,  ombre = true, barre=none, logo=\bcbook, marge=4]{~\textbf{Théorème} \textbf{\thethme} {\itshape #1} }   \par}
{
\end{bclogo}
 \par \bigskip}

 \newenvironment{theoremedef}[1]
{\par \medskip   \addtocounter{thme}{1} \noindent  
\begin{bclogo}[arrondi =0.1,  ombre = true, barre=none, logo=\bcbook, marge=4]{~\textbf{Théorème-Définition} \textbf{\thethme} {\itshape #1} }   \par}
{
\end{bclogo}
 \par \bigskip }
 
\newcounter{prop}
\newenvironment{propriete}[1]
{\par \medskip   \addtocounter{prop}{1} \noindent  
\begin{bclogo}[arrondi =0.1,  ombre = true, barre=none, logo=\bcbook, marge=4]{~\textbf{Propriété} \textbf{\theprop} {\itshape #1} }   \par}
{
\end{bclogo}
 \par \bigskip }


\newenvironment{corollaire}[1]
{\par \medskip   \noindent  
\begin{bclogo}[arrondi =0.1,  ombre = true, barre=none, logo=\bcbook, marge=4]{~\textbf{Corollaire} {\itshape #1} } \par }
{
\end{bclogo}
 \par \bigskip }

\newenvironment{demo}[1]
{\par \medskip   \noindent  
\begin{bclogo}[arrondi =0.1,  ombre = true, barre=zigzag, noborder = true, logo=\bcloupe, marge=0]{~\textbf{Démonstration} {\itshape #1} } \par \vspace{10pt}}
{
\end{bclogo}
 \par \bigskip }

\newcounter{activite}
\newenvironment{activite}[1]
{\par \medskip   \noindent   \addtocounter{activite}{1}
\begin{bclogo}[arrondi =0.1,   noborder = true, logo=\bcvelo, marge=4]{~\textbf{Activité} \textbf{\theactivite} {\itshape #1} }  \par}
{
\end{bclogo}
 \par \bigskip }


\newenvironment{synthese}
{\par \medskip   \noindent   
\begin{bclogo}[arrondi =0.1,   noborder = true, logo=\bccle, marge=4]{~\textbf{Synthèse}   }  \par}
{
\end{bclogo}
 \par \bigskip }
 
 
\newcounter{rque}
\newenvironment{remarque}
{\par \medskip    \addtocounter{rque}{1} \noindent  
\begin{bclogo}[arrondi =0.1,  ombre = true, barre=snake, noborder = true, logo=\bcinfo, marge=0]{~\textbf{Remarque} \textbf{\therque}}  \par }
{
\end{bclogo}
 \par \bigskip }

\newcounter{def}
\newenvironment{definition}[1]
{\par \medskip   \addtocounter{def}{1} \noindent  
\begin{bclogo}[arrondi =0.1,  ombre = true, barre=none, logo=\bcbook, marge=4]{~\textbf{Définition} \textbf{\thedef} {\itshape #1} }  \par}
{
\end{bclogo}
 \par \bigskip }
 
 
 \newcounter{cours}
\newenvironment{cours}[1]
{\par \medskip   \addtocounter{cours}{1} \noindent  
\begin{bclogo}[arrondi =0.1,  ombre = true, barre=none, logo=\bcbook, marge=4]{~\textbf{Point de cours} \textbf{\thecours} {\itshape #1} }  \par}
{
\end{bclogo}
 \par \bigskip }
 
 

\newenvironment{introduction}
{\par \medskip    \noindent  
 \begin {bclogo}[couleur = blue!5 , arrondi =0.1,logo=\bcrosevents, marge=4] {~\textbf{Introduction}    }
 \par }
{
\end{bclogo}
 \par \bigskip }
 
\newenvironment{memo}[1]
{\par \medskip    \noindent  
\begin{bclogo}[arrondi =0.1,  ombre = true, barre=none, logo=\bccle, marge=4]{~\textbf{À retenir}  {\itshape #1} }  \par}
{
\end{bclogo}
 \par \bigskip }
 
\newcounter{exple}
\newenvironment{exemple}[1]
{\par \medskip   \addtocounter{exple}{1} \noindent  
\begin{bclogo}[arrondi =0.1,   noborder = true, logo=\bclampe, marge=4]{~\textbf{Exemple} \textbf{\theexple} {\itshape #1} }  \par}
{
\end{bclogo}
 \par \bigskip }




\newcounter{alg}
\newenvironment{algorithme}[1]
{\par \medskip   \addtocounter{alg}{1} \noindent  
 \begin {bclogo}[noborder = true, barre=zigzag,logo=\bcpython, marge=4] {~\textbf{Algorithmique} \textbf{\thealg} {\itshape #1} }  \par}
{
\end{bclogo}
 \par \bigskip }

\newcounter{prog}
\newenvironment{programme}[1]
{\par \medskip   \addtocounter{prog}{1} \noindent  
 \begin {bclogo}[noborder = true, barre=zigzag,logo=\bcpython, marge=4] {~\textbf{Programme} \textbf{\theprog} {\itshape #1} }  \par  \bigskip}
{
\end{bclogo}
 \par \bigskip }
 
\newcounter{logi}
\newenvironment{logique}[1]
{\par \medskip   \addtocounter{logi}{1} \noindent  
 \begin {bclogo}[noborder = true, barre=zigzag,logo=\bclampe, marge=4] {~\textbf{Logique} \textbf{\thelogi} {\itshape #1} }  \par}
{
\end{bclogo}
 \par \bigskip }


\newenvironment{methode}[1]
{\par \medskip    \noindent  
 \begin {bclogo}[arrondi =0.1,logo=\bcoutil, marge=4,noborder = true] {~\textbf{Méthode}   {\itshape #1} }  \par}
{
\end{bclogo}
 \par \bigskip }


\newcounter{histo}
\newenvironment{histoire}[1]
{\par \medskip   \addtocounter{histo}{1} \noindent  
 \begin {bclogo}[couleur = blue!10 , arrondi =0.1,logo=\bchorloge, marge=4] {~\textbf{Histoire} \textbf{\thehisto} {\itshape #1} }  \par}
{
\end{bclogo}
 \par \bigskip }




%Environnement contenu pour un document présentant une progression annuelle
\newenvironment{contenu}
{\par \medskip   \begin {bclogo}[ noborder = true,logo=\bccrayon] \noindent {\large \textbf{Contenu de la séance}} \vspace*{10pt} \par  }
{\end{bclogo}  \par \medskip }

%Environnement programme pour un document présentant une progression annuelle
%\newenvironment{programme}
%{\par \medskip   \begin {bclogo}[ noborder = true, barre=zigzag,logo=\bcinfo] \noindent {\large \textbf{Programme officiel}} \vspace*{10pt} \par  }
%{\end{bclogo}  \par \medskip }

%Environnement programme pour un document présentant une progression annuelle
\newenvironment{ressource}
{\par \medskip   \begin {bclogo}[ noborder = true,logo=\bcbook] \noindent {\large \textbf{Ressources}}\\vspace*{10pt} \par }
{\end{bclogo}  \par \medskip }




%%%%%%%%%%%%%%%%%%Maths divers%%%%%%%%%%%%%%%%%%%%%%%%%
%%%%%%%%%%%%%Nombres%%%%%%%%%%%%%%%%

%Ensemble prive de...
%\newcommand{\prive}{\boi}%{\backslash}

%Ensembles de nombres%%%%%%%%%%%%%%%%%
\newcommand{\R}{\mathbb{R}}
\newcommand{\N}{\mathbb{N}}
\newcommand{\D}{\mathbb{D}}
\newcommand{\Z}{\mathbb{Z}}
\newcommand{\Q}{\mathbb{Q}}
%\newcommand{\C}{\mathbb{C}}
\newcommand{\df}{~\ensuremath{]0;+\infty[}~}
\newcommand{\K}{\mathbb{K}}

%%%%%%%%Arithmetique%%%%%%%%%%
%PGCD, PPCM
\newcommand{\PGCD}{\mathop{\rm PGCD}\nolimits}
\newcommand{\PPCM}{\mathop{\rm PPCM}\nolimits}

%Intervalles
\newcommand{\interoo}[2]{]#1\, ;\, #2[}
\newcommand{\Interoo}[2]{\left]#1\, ;\, #2\right[}
\newcommand{\interof}[2]{]#1\, ;\, #2]}
\newcommand{\Interof}[2]{\left]#1\, ;\, #2\right]}
\newcommand{\interfo}[2]{[#1\, ;\, #2[}
\newcommand{\Interfo}[2]{\left[#1\, ;\, #2\right[}
\newcommand{\interff}[2]{[#1\, ;\, #2]}
\newcommand{\Interff}[2]{\left[#1\, ;\, #2\right]}
%\newcommand\interentiers #1#2{[\! [#1\, ;\, #2]\! ]}
\newcommand{\interentiers}[2]{\llbracket #1\, ;\, #2\rrbracket}
%


%%%%%%%%%%%%%%Nombres complexes%%%%%

\newcommand{\ic}{\text{i}}
%\newcommand{\I}{\text{i}}
\newcommand{\im}[1]{\text{Im}\left(#1\right)}
\newcommand{\re}[1]{\text{Re}\left(#1\right)}
\newcommand{\Arg}[1]{\text{arg}\left(#1\right)}
\newcommand{\Mod}[1]{\left[#1\right]}
%Parties entière, réelle, imaginaire, nombre i
\newcommand{\ent}[1]{\text{E}\left(#1\right)}
\renewcommand{\Re}{\mathop{\rm Re}\nolimits}
\renewcommand{\Im}{\mathop{\rm Im}\nolimits}
\renewcommand{\i}{\textrm{i}}

%%%%%%%%%%%Probabilites et statistiques%%%%%
\newcommand{\loibinom}[2]{\mathcal{B}\left(#1\ ; \ #2 \right)}
\newcommand{\loinorm}[2]{\mathcal{N}\left(#1\ ; \ #2 \right)}
\newcommand{\loiexp}[1]{\mathcal{E}\left(#1\right)}
\newcommand{\proba}[1]{\mathbb{P}\big(#1\big)}
\newcommand{\probacond}[2]{\mathbb{P}_{#2}\big(#1\big)}
\newcommand{\esperance}[1]{\mathbb{E}\left(#1\right)}
\newcommand{\variance}[1]{\mathbb{V}\left(#1\right)}
\newcommand{\ecart}[1]{\sigma\left(#1\right)}
\newcommand{\dnormx}{\frac{1}{\sqrt{2\pi}} \text{e}^{-\frac{x^2}{2}}}
\newcommand{\dnormt}{\frac{1}{\sqrt{2\pi}} \text{e}^{-\frac{t^2}{2}}}

%Covariance
\newcommand{\cov}{\mathop{\rm cov}\nolimits}
%


%%%%%%%%%%Analyse%%%%%%%%%%%

%%%%%%%%%%%Courbe%%%%%%%%%%%%
\newcommand{\courbe}[1]{\ensuremath{\mathcal{C}_{#1}}}

%%%%%%%Fonction exponentielle%%%%%
\newcommand{\fe}{~fonction exponentielle~}
\newcommand{\e}{\text{e}}

%Fonction cotangente
\newcommand{\cotan}{\mathop{\rm cotan}\nolimits}
%%%%%%%%%%%%%%%%%%%%%%%%%%%%%%%%%%%%%%%%%
%
%Fonctions hyperboliques
\newcommand{\ch}{\mathop{\rm ch}\nolimits}
\newcommand{\sh}{\mathop{\rm sh}\nolimits}


%%%%%%%%%%%%%%Limites%%%%%%
\newcommand{\limite}[2]{\lim\limits
_{x \to #1} #2}
\newcommand{\limitesuite}[1]{\lim\limits
_{n \to +\infty} #1}
\newcommand{\limiteg}[2]{\lim\limits
_{\substack{x \to #1 \\ x < #1 }} #2}
\newcommand{\limited}[2]{\lim\limits
_{\substack{x \to #1 \\ x > #1 }} #2}

%%%%%%%%%%Continuité%%%%%%%%%%%
\newcommand{\TVI}{théorème des valeurs intermédiaires}

%%%%%%%%%%%Suites%%%%%%%%%%%%
\newcommand{\suite}[1]{\ensuremath{\left(#1_{n}\right)}}
\newcommand{\Suite}[2]{\ensuremath{\left(#1\right)_{#2}}}
%

%%%%%%%%%%%%%%%Calcul intégral%%%%%%
\newcommand{\dx}{\ensuremath{\text{d}x}}		% dx
\newcommand{\dt}{\ensuremath{\text{d}t}}		% dt
\newcommand{\dtheta}{\ensuremath{\text{d}\theta}}		% dtheta
\newcommand{\dy}{\ensuremath{\text{d}y}}		% dy
\newcommand{\dq}{\ensuremath{\text{d}q}}		% dq

%%%Intégrale%%%
\newcommand{\integralex}[3]{\int_{#1}^{#2} #3 \ \dx}
\newcommand{\integralet}[3]{\int_{#1}^{#2} #3 \ \dt}
\newcommand{\integraletheta}[3]{\int_{#1}^{#2} #3 \ \dtheta}

%%%%%Equivalent%%
\newcommand{\equivalent}[1]{\build\sim_{#1}^{}}

%o et O%%%%
\renewcommand{\o}[2]{\build o_{#1\to #2}^{}}
\renewcommand{\O}[2]{\build O_{#1\to #2}^{}}



%%%%%%%%%%%%%%%Geometrie%%%%%%%%%%%%%%%%%%%%%%%

%%%%%%%%%%%%%%%Reperes%%%%%%%%%%%%%%
\def\Oij{\ensuremath{\left(\text{O},~\vect{\imath},~\vect{\jmath}\right)}}
\def\Oijk{\ensuremath{\left(\text{O},~\vect{\imath},~ \vect{\jmath},~ \vect{k}\right)}}
\def\Ouv{\ensuremath{\left(\text{O},~\vect{u},~\vect{v}\right)}}
\renewcommand{\ij}{(\vec\imath\, ;\vec\jmath\,)}
\newcommand{\ijk}{(\vec\imath\, ;\vec\jmath\, ;\vec k\,)}
\newcommand{\OIJ}{(O\,;\, I\,;\, J\,)}
\newcommand{\repere}[3]{\big(#1\, ;\,\vect{#2} ;\vect{#3}\big)}
\newcommand{\reperesp}[4]{\big(#1\, ;\,\vect{#2} ;\vect{#3} ;\vect{#4}\big)}

%%%%%%%%%Coordonnees%%%%%%%%%%%%%%
\newcommand{\coord}[2]{(#1\, ;\, #2)}
\newcommand{\bigcoord}[2]{\big(#1\, ;\, #2\big)}
\newcommand{\Coord}[2]{\left(#1\, ;\, #2\right)}
\newcommand{\coordesp}[3]{(#1\, ;\, #2\, ;\, #3)}
\newcommand{\bigcoordesp}[3]{\big(#1\, ;\, #2\, ;\, #3\big)}
\newcommand{\Coordesp}[3]{\left(#1\, ;\, #2\, ;\, #3\right)}
\newcommand{\Vcoord}[3]{\begin{pmatrix} #1 \\ #2 \\ #3 \end{pmatrix}}
%Symboles entre droites
%\newcommand{\paral}{\sslash}
\newcommand{\paral}{\mathop{/\!\! /}}
%

%%%%%%%%%Produit scalaire, Angles%%%%%%%%%%
\newcommand{\scal}[2]{\vect{#1} \, \cdot \, \vect{#2}}
\newcommand{\Angle}[2]{\left(\vect{#1} \, , \, \vect{#2}\right)}
\newcommand{\Anglegeo}[2]{\left(\widehat{\vect{#1} \, ; \, \vect{#2}}\right)}
\renewcommand{\angle}[1]{\widehat{#1}}
\newcommand{\anglevec}[2]{\left(\vect {#1}\, ,\,\vect {#2} \right)}
\newcommand{\anglevecteur}[2]{(#1\, , \, #2)}
\newcommand{\Anglevec}[2]{(\vecteur{#1}\, ,\,\vecteur{#2})}
\newcommand{\prodscal}[2]{#1 \, \cdot \, #2}
%


%Arc
%\newcommand{\arc}[1]{\wideparen{#1}}
\newcommand{\arcoriente}[1]{\overset{\curvearrowright}{#1}}
%
%


%%%%%%%%%%%%%%%Normes%%%%%%%%%%%%%%%%
\newcommand{\norme}[1]{\left\| #1\right\|}
\newcommand{\normebis}[1]{\delim{2pt}{\|}{9pt}\! #1\delim{2pt}{\|}{9pt}}
\newcommand{\normetriple}[1]{\left |\kern -.07em\left\| #1\right |\kern -.07em\right\|}
\newcommand{\valabs}[1]{\big| \, #1 \, \big|}
%

%%%%%%%%%%%%%%%%%%%%%%%%%%%Degré%%%%%%
%\newcommand{\Degre}{\ensuremath{^\circ}}
%La commande \degre est déjà définie dans le package babel

%%%%%%%%%%Vecteurs%%%%%%%%%%%
\newcommand{\vect}[1]{\mathchoice%
{\overrightarrow{\displaystyle\mathstrut#1\,\,}}%
{\overrightarrow{\textstyle\mathstrut#1\,\,}}%
{\overrightarrow{\scriptstyle\mathstrut#1\,\,}}%
{\overrightarrow{\scriptscriptstyle\mathstrut#1\,\,}}}



%%%%%%%%%%%%%Algebre%%%%%%%%%%%%%%%


%%%%%%%%%%Systemes%%%%%%%%%%%
%Systemes
\newcommand{\sys}[2]{
\left\lbrace
 \begin{array}{l}
  \negthickspace\negthickspace #1\\
  \negthickspace\negthickspace #2\\
 \end{array}
\right.\negthickspace\negthickspace}
\newcommand{\Sys}[3]{
\left\lbrace
 \begin{array}{l}
  #1\\
  #2\\
  #3\\
 \end{array}
\right.}
\newcommand{\Sysq}[4]{
\left\lbrace
 \begin{array}{l}
  #1\\
  #2\\
  #3\\
  #4\\
 \end{array}
\right.}
%
%

%%%%%%%%%%%%%%%%Matrices%%%%%%%%%%%%%%%%%%
%Comatrice
\newcommand{\com}{\mathop{\rm com}\nolimits}
%
%
%Trace
\newcommand{\tr}{\mathop{\rm tr}\nolimits}
%
%
%Transposee
\newcommand{\transposee}[1]{{\vphantom{#1}}^t\negmedspace #1}
%
%
%Noyau
\newcommand{\Ker}{\mathop{\rm Ker}\nolimits}
%
%

%
%Matrices
\newcommand{\Mn}{\mathcal M_n}
\newcommand{\matrice}[4]{
\left(
 \begin{array}{cc}
  #1 & #2 \\
  #3 & #4
 \end{array}
\right)}

\newcommand{\Matrice}[9]{
\left(
 \begin{array}{ccc}
  #1 & #2 & #3\\
  #4 & #5 & #6\\
  #7 & #8 & #9
 \end{array}
\right)}
\newcommand{\Vect}[3]{
\left(\negmedspace
 \begin{array}{c}
  #1\\
  #2\\
  #3
 \end{array}\negmedspace
\right)}
\newcommand{\Ideux}{\matrice{1}{0}{0}{1}}
\newcommand{\Itrois}{\Matrice{1}{0}{0}{0}{1}{0}{0}{0}{1}}
%
%
%Determinants
\newcommand{\determinant}[4]{
\left|
 \begin{array}{cc}
  #1 & #2 \\
  #3 & #4
 \end{array}
\right|}
\newcommand{\Determinant}[9]{
\left|
 \begin{array}{ccc}
  #1 & #2 & #3\\
  #4 & #5 & #6\\
  #7 & #8 & #9
 \end{array}
\right|}

\begin{document}
\maketitle

\renewcommand*\contentsname{Table des matières}
{
\hypersetup{linkcolor=}
\setcounter{tocdepth}{3}
\tableofcontents
}
\hypertarget{cruxe9dits}{%
\section*{Crédits}\label{cruxe9dits}}
\addcontentsline{toc}{section}{Crédits}

\emph{Ce cours est largement inspiré du chapitre 29 du manuel NSI de la
collection Tortue chez Ellipse, auteurs : Ballabonski, Conchon,
Filliatre, N'Guyen. J'ai également consulté le prepabac Première NSI de
Guillaume Connan chez Hatier, la documentation en ligne de la fondation
Mozilla \url{https://developer.mozilla.org/fr/docs/Apprendre/JavaScript}
et le tutoriel de w3schools
\url{https://www.w3schools.com/js/default.asp}.}

\hypertarget{javascriptjavascript-et-la-programmation-uxe9vuxe9nementielle-cuxf4tuxe9-client}{%
\section{\texorpdfstring{\href{https://developer.mozilla.org/fr/docs/Glossaire/JavaScript}{Javascript}
et la programmation événementielle côté
client}{Javascript et la programmation événementielle côté client}}\label{javascriptjavascript-et-la-programmation-uxe9vuxe9nementielle-cuxf4tuxe9-client}}

\begin{cours}{}

Dans le chapitre précédent, on a présenté des exemples de pages Web
\textbf{dynamiques} générées par des programmes en
\href{https://developer.mozilla.org/fr/docs/Glossaire/PHP}{PHP} ou
\href{https://docs.python.org/3.7/library/cgi.html}{Python}. Chaque mise
à jour de la page nécessite donc un nouveau cycle requête/réponse entre
le client et le serveur. C'est indispensable s'il s'agit de modifier
l'état d'une ressource côté serveur (une base de données par exemple),
mais les changements peuvent n'être que temporaires et n'affecter que
des éléments de la page côté client. C'est le cas de l'exemple déjà
traité en
\href{https://developer.mozilla.org/fr/docs/Glossaire/PHP}{PHP} de la
conversion d'unité pour une mesure de température.

\href{https://developer.mozilla.org/fr/docs/Glossaire/JavaScript}{Javascript}
est un langage interprété qui répond à ce besoin qui s'exécute dans le
navigateur du client.
\href{https://developer.mozilla.org/fr/docs/Glossaire/JavaScript}{Javascript}
s'est imposé depuis son apparition en 1995 dans le navigateur
\href{https://fr.wikipedia.org/wiki/Netscape_Navigator}{Netscape} comme
le principal langage de développement Web en \emph{frontend} (côté
client) et depuis une dizaine d'années, sa variante
\href{https://developer.mozilla.org/fr/docs/Glossaire/Node.js}{Node.js}
concurence les langages de développement \emph{backend} (côté serveur)
comme \href{https://developer.mozilla.org/fr/docs/Glossaire/PHP}{PHP} ou
\href{https://docs.python.org/3.7/library/cgi.html}{Python}.

Une page Web moderne, reçue par un client, comporte au moins trois
composants logiciels :

\begin{itemize}
\tightlist
\item
  \href{https://developer.mozilla.org/fr/docs/Glossaire/HTML}{HTML} pour
  la structure du document.
\item
  \href{https://developer.mozilla.org/fr/docs/Glossaire/CSS}{CSS} pour
  le paramétrage de l'apparence des éléments et leur positionnement.
\item
  \href{https://developer.mozilla.org/fr/docs/Glossaire/JavaScript}{Javascript}
  pour la définition de programmes qui vont réagir à des événements
  déclenchés par l'utilisateur de la page et modifier la structure de
  données de la page ( éléments
  \href{https://developer.mozilla.org/fr/docs/Glossaire/HTML}{HTML} et
  styles
  \href{https://developer.mozilla.org/fr/docs/Glossaire/CSS}{CSS}) à
  travers
  l'\href{https://developer.mozilla.org/fr/docs/Glossaire/API}{API}
  \footnote{Note :
    \href{https://developer.mozilla.org/fr/docs/Glossaire/API}{API} est
    l'acronyme d'Application Programming Interface} nommée
  \href{https://developer.mozilla.org/fr/docs/Glossaire/DOM}{DOM}
  \footnote{Note :
    \href{https://developer.mozilla.org/fr/docs/Glossaire/DOM}{DOM} est
    l'acronyme de Document Object Model}. Le
  \href{https://developer.mozilla.org/fr/docs/Glossaire/DOM}{DOM} est
  une représentation de l'ensemble de la structure de la page Web sous
  la forme d'un arbre : les noeuds sont les éléments
  \href{https://developer.mozilla.org/fr/docs/Glossaire/HTML}{HTML} et
  ils ont une liste de propriétés : contenu, style, événements associés
  \ldots{} L'inspecteur des outils de développement, accessibles avec la
  touche \passthrough{\lstinline!F12!} dans un navigateur, permet de
  visualiser et modifier les propriétés des éléments du
  \href{https://developer.mozilla.org/fr/docs/Glossaire/DOM}{DOM}.
\end{itemize}

L'environnement d'exécution d'un code
\href{https://developer.mozilla.org/fr/docs/Glossaire/JavaScript}{Javascript}
est confiné à l'onglet de la page Web où il est chargé. Pour des raisons
de sécurité il ne doit pas interagir avec d'autres pages ou des
ressources du poste client. Par ailleurs, si on recharge la page auprès
du serveur, l'environnement
\href{https://developer.mozilla.org/fr/docs/Glossaire/JavaScript}{Javascript}
est réinitialisé et les modifications de la page effectuées par un code
\href{https://developer.mozilla.org/fr/docs/Glossaire/JavaScript}{Javascript}
ne sont pas répercutées sur la page source disponible sur le serveur.

Le schéma ci-dessous illustre le fonctionnement du
\href{https://developer.mozilla.org/fr/docs/Glossaire/JavaScript}{Javascript}
qui correspond à un paradigme de
\href{https://fr.wikipedia.org/wiki/Programmation_\%C3\%A9v\%C3\%A9nementielle}{programmation
événementielle}.

\end{cours}

\begin{figure}
\centering
\includegraphics{images/schema-dom.png}
\caption{``Javascript et DOM''}
\end{figure}

\hypertarget{premiers-pas-dans-la-console-javascriptjavascript}{%
\section{\texorpdfstring{Premiers pas dans la console
\href{https://developer.mozilla.org/fr/docs/Glossaire/JavaScript}{Javascript}}{Premiers pas dans la console Javascript}}\label{premiers-pas-dans-la-console-javascriptjavascript}}

\begin{cours}{}

Compléter ce tableau sur les opérateurs en
\href{https://developer.mozilla.org/fr/docs/Glossaire/JavaScript}{Javascript}
à partir de la page
\url{https://developer.mozilla.org/fr/docs/Web/JavaScript/Guide/Expressions_et_Op\%C3\%A9rateurs}.

\begin{longtable}[]{@{}cc@{}}
\toprule
Opérateurs & Description\tabularnewline
\midrule
\endhead
\passthrough{\lstinline!=!} &\tabularnewline
\passthrough{\lstinline!*!} &\tabularnewline
\passthrough{\lstinline!/!} &\tabularnewline
\passthrough{\lstinline!**!} &\tabularnewline
\passthrough{\lstinline!==!} ou \passthrough{\lstinline!===!}
&\tabularnewline
\passthrough{\lstinline"!="} ou \passthrough{\lstinline"!==="}
&\tabularnewline
\passthrough{\lstinline!\&\&!} &\tabularnewline
\passthrough{\lstinline!||!} &\tabularnewline
\passthrough{\lstinline"!"} &\tabularnewline
\bottomrule
\end{longtable}

\end{cours}

\begin{exercice}{}

Ouvrir un nouvel onglet dans un navigateur Web.

\begin{enumerate}
\def\labelenumi{\arabic{enumi}.}
\item
  Ouvrir la console
  \href{https://developer.mozilla.org/fr/docs/Glossaire/JavaScript}{Javascript}
  dans la fenêtre des outils de développement avec
  \passthrough{\lstinline!F12!} ou
  \passthrough{\lstinline!CTRL + SHIFT + K!} sous Firefox. On va
  exécuter de façon interactive, une séquence d'instructions
  \href{https://developer.mozilla.org/fr/docs/Glossaire/JavaScript}{Javascript}
  pour présenter quelques constructions élémentaires et propriétés du
  langage. Chaque instruction pourra modifier l'état du
  \href{https://developer.mozilla.org/fr/docs/Glossaire/DOM}{DOM} et
  donc le rendu graphique de la page Web.
\item
  On commence par quelques manipulations de variables et calculs :
\end{enumerate}

\begin{lstlisting}
>>> let a = 1
undefined
>>> (a * 3 + 1) ** 2 / 5 - 1
2.2
  >>> (a * 3 + 1) ** 2 // commentaire !
16
>>> let b = "Hello"
undefined
>>> b + " World"
"Hello World"
>>> typeof(a)
"number"
>>> typeof(b)
"string"
>>> a = b + a
"Hello1"
>> typeof(a)
"string"
\end{lstlisting}

\begin{itemize}
\tightlist
\item
  Barrer les propositions fausses : en
  \href{https://developer.mozilla.org/fr/docs/Glossaire/JavaScript}{Javascript}
  est à typage \emph{(fort \textbar{} faible)} et \emph{(dynamique
  \textbar{} statique)} et une variable égale à 5 se déclare avec
  \emph{(let a = 5 \textbar{} a = 5)}.
\end{itemize}

\begin{enumerate}
\def\labelenumi{\arabic{enumi}.}
\setcounter{enumi}{2}
\tightlist
\item
  Examinons un exemple avec une fonction, une structure conditionnelle
  et une boucle. Dans la console, passer en mode éditeur multiligne avec
  \passthrough{\lstinline!CTRL + B!} et saisir le code ci-dessous :
\end{enumerate}

\begin{lstlisting}
function valabs(x){
  if (x < 0){
    return -x;
  }
  else {
    return x;
  }
}

for (let i = -4; i < 5; i = i + 1){
  if (valabs(i) > 2){
    alert(i);
  }
  else {
    console.log(i);
  }
}
\end{lstlisting}

\begin{itemize}
\tightlist
\item
  Barrer les propositions fausses : en
  \href{https://developer.mozilla.org/fr/docs/Glossaire/JavaScript}{Javascript},
  les blocs d'instructions sont délimités par \emph{(l'indentation
  \textbar{} des accolades)}, les fonctions sont déclarées avec le mot
  clef \emph{(def \textbar{} function)} et une boucle inconditionnelle
  sur les entiers entre 1 et 10 commence par l'instruction \emph{( for k
  in range(1, 11) \textbar{} for (let k = 1; k \textless{} 11; k = k +
  1) )}
\item
  Barrer les propositions fausses : la fonction
  \passthrough{\lstinline!alert!} affiche son paramètre dans \emph{(une
  fenêtre pop-up \textbar{} la console)} tandis que la fonction
  \passthrough{\lstinline!console.log!} affiche son paramètre dans
  \emph{(une fenêtre pop-up \textbar{} la console)}.
\end{itemize}

\end{exercice}

\begin{exercice}{}

\href{https://developer.mozilla.org/fr/docs/Glossaire/JavaScript}{Javascript}
dispose de toutes les constructions permettant de progammer les mêmes
algorithmes qu'en
\href{https://docs.python.org/3.7/library/cgi.html}{Python}, mais s'il a
été inventé par les développeurs de
\href{https://fr.wikipedia.org/wiki/Netscape_Navigator}{Netscape} c'est
pour manipuler les éléments
\href{https://developer.mozilla.org/fr/docs/Glossaire/HTML}{HTML}, à
travers l'interface du
\href{https://developer.mozilla.org/fr/docs/Glossaire/DOM}{DOM}. Il
existe plusieurs façons de capturer un élément
\href{https://developer.mozilla.org/fr/docs/Glossaire/HTML}{HTML}, la
plus simple s'il s'agit d'un élément particulier est
\href{https://www.w3schools.com/js/js_htmldom_methods.asp}{l'accès par
son identifiant unique \passthrough{\lstinline!id!}} (à condition qu'il
soit défini).\\
Ouvrir un navigateur Web et charger la page
\url{https://frederic-junier.org/NSI/sandbox/bac.html}. Le code source
\href{https://developer.mozilla.org/fr/docs/Glossaire/HTML}{HTML} de la
page est donné ci-dessous :

\begin{lstlisting}[language=HTML]
<!DOCTYPE html>
<head>
  <title> "Bac à sable"  </title>
   <meta charset="utf-8">
<body>
<ul id="langages">
    <li> HTML </li>
    <li> CSS </li>
</ul>
</body>
</html>
\end{lstlisting}

\begin{enumerate}
\def\labelenumi{\arabic{enumi}.}
\tightlist
\item
  Saisir successivement les instructions ci-dessous dans la console,
  observer ce qu'il se passe dans la page étape par étape. Caractériser
  le rôle de chaque instruction puis saisir des instructions qui
  permettent de modifier le contenu
  \href{https://developer.mozilla.org/fr/docs/Glossaire/HTML}{HTML} et
  le style du second item de la liste.
\end{enumerate}

\begin{lstlisting}
>>> let list = document.getElementById("langages")
undefined
>>> list.style.color = "red"
"red"
>>> let item1 = list.children[0]
undefined
>>> item1
<li>
>>> item1.innerHTML = 'HTML5'
"HTML5"
\end{lstlisting}

\includegraphics{images/exo2-fig1.png}\\

\begin{enumerate}
\def\labelenumi{\arabic{enumi}.}
\setcounter{enumi}{1}
\tightlist
\item
  Saisir successivement les instructions ci-dessous dans la console,
  observer ce qu'il se passe dans la page étape par étape. Caractériser
  le rôle de chaque instruction puis saisir des instructions qui
  permettent de rajouter un quatrième item à la liste dont le contenu
  est \href{https://fr.wikipedia.org/wiki/Kotlin_(langage)}{Kotlin}.
\end{enumerate}

\begin{lstlisting}
>>> let item3 = document.createElement("li")
undefined
>>> list.appendChild(item3)
<li>
>>> item3.innerHTML = "Javascript"
"Javascript"
\end{lstlisting}

\begin{enumerate}
\def\labelenumi{\arabic{enumi}.}
\setcounter{enumi}{2}
\tightlist
\item
  Passer la console en mode mutliligne avec
  \passthrough{\lstinline!CTRL + B!} puis saisir et exécuter le code
  ci-dessous. Survoler la liste ou cliquer dessus avec la souris et
  observer ce qui se passe dans la page. Caractériser le rôle de chaque
  instruction en vous appuyant sur la lecture de l'article
  \url{https://developer.mozilla.org/fr/docs/Web/API/GlobalEventHandlers/onclick}.
\end{enumerate}

\begin{lstlisting}
function apparition(){
  list.style.visibility = "visible";
}

function disparition(){
  list.style.visibility = "hidden";
}

list.onmouseover = disparition ;  
//le symbole ; est le séparateur d'instruction en Javascript
list.onmouseleave = apparition ;
list.onclick = function() { alert("Clic !") ; } ;  //définition d'une fonction anonyme
\end{lstlisting}

\begin{enumerate}
\def\labelenumi{\arabic{enumi}.}
\setcounter{enumi}{3}
\tightlist
\item
  Recharger la page avec \passthrough{\lstinline!F5!}. L'environnement
  \href{https://developer.mozilla.org/fr/docs/Glossaire/JavaScript}{Javascript}
  a-t-il été conservé ?
\end{enumerate}

\end{exercice}

\hypertarget{approfondissement-ruxe9daction-de-scripts-et-gestionnaires-duxe9vuxe9nements}{%
\section{Approfondissement : rédaction de scripts et gestionnaires
d'événements}\label{approfondissement-ruxe9daction-de-scripts-et-gestionnaires-duxe9vuxe9nements}}

\begin{exemple}{}

Ouvrir un navigateur Web et accéder à la page
\url{https://frederic-junier.org/NSI/sandbox/compteur-clic.html}.

On donne ci-dessous le code source de la page Web.

\begin{lstlisting}[language=HTML]
<!DOCTYPE html>

<html lang="fr">

<head><!-- Début en-tête -->
  <title>Compteur de clics </title>
  <meta charset="utf-8">
  <!-- Début script JS --> 
  <script>  

      let compteur = 0;

      function suivant(){
          compteur = compteur + 1;
          let v = document.getElementById("valeur");
          if (compteur <= 1) {
            v.innerHTML = compteur + " clic";
          }
          else {
            v.innerHTML = compteur + " clics";
          }          
      }

      function reinitialiser(){
          compteur = 0;
          let v = document.getElementById("valeur");
          v.innerHTML = compteur + " clic";
      }

  </script>  <!-- Fin script JS -->
</head><!-- Fin en-tête -->
 
<body><!-- Début corps -->   
  <h1>Compteur de clics en Javascript </h1>
    <button onclick="suivant();">Clic !</button>
    <br>
    <span id="valeur">0</span>
    <br>
    <button onclick="reinitialiser();">Réinitialiser !</button>
</body><!-- Fin corps -->

</html> 
\end{lstlisting}

\begin{enumerate}
\def\labelenumi{\arabic{enumi}.}
\item
  On peut interagir avec la page en cliquant sur le bouton \emph{Clic !}
  puis réinitialiser le compteur en cliquant sur le bouton
  \emph{Réinitialiser}.
\item
  Ouvrir l'onglet inspecteur dans la fenêtre des outils de développement
  et afficher la structure de la page comme ci-dessous.

  \begin{itemize}
  \tightlist
  \item
    Repérer l'étiquette \passthrough{\lstinline!event!} attachée au
    bouton de contenu \emph{Clic !} et déplier son contenu.
  \end{itemize}

  \includegraphics{images/exo3-fig1.png}~

  \begin{itemize}
  \tightlist
  \item
    Pour cet élément \passthrough{\lstinline!button!}, l'attribut
    \passthrough{\lstinline!onclick="suivant()"!} a défini une fonction
    \href{https://developer.mozilla.org/fr/docs/Glossaire/JavaScript}{Javascript}
    \passthrough{\lstinline!function onclick(event) \{suivant();\}!} qui
    prend en paramètre un \textbf{événement utilisateur} (ici un click)
    et réagit en appelant la fonction \passthrough{\lstinline!suivant!}.
    La fonction \passthrough{\lstinline!suivant!} est appelée
    \textbf{gestionnaire d'événement}.
  \item
    Le code de la fonction \passthrough{\lstinline!suivant!} se trouve
    dans la balise \passthrough{\lstinline!<script>!}. Placée dans
    l'en-tête, celle-ci rassemble tout le code
    \href{https://developer.mozilla.org/fr/docs/Glossaire/JavaScript}{Javascript}
    chargé dans la page.
  \item
    On a vu précédemment que la fonction est appelée dès qu'un événement
    \passthrough{\lstinline!click!} se produit sur le premier élément
    \passthrough{\lstinline!<button>!}. En examinant le code de
    \passthrough{\lstinline!suivant!}, on voit qu'un lien est d'abord
    créé entre l'élément \passthrough{\lstinline!<span>!} d'identifiant
    \passthrough{\lstinline!"valeur"!} et une variable
    \href{https://developer.mozilla.org/fr/docs/Glossaire/JavaScript}{Javascript}
    à l'aide d'une méthode
    \href{https://developer.mozilla.org/fr/docs/Glossaire/DOM}{DOM} par
    l'instruction :
  \end{itemize}

\begin{lstlisting}
let v = document.getElementById("valeur");
\end{lstlisting}

  \begin{itemize}
  \tightlist
  \item
    Ensuite on peut manipuler l'élément
    \href{https://developer.mozilla.org/fr/docs/Glossaire/HTML}{HTML} et
    modifier son contenu avec :
  \end{itemize}

\begin{lstlisting}
v.innerHTML = compteur + " clic";
\end{lstlisting}

  \begin{itemize}
  \item
    La variable \passthrough{\lstinline!compteur!} est définie en dehors
    de \passthrough{\lstinline!suivant!} car c'est une variable globale.
    Les variables globales sont essentielles en
    \href{https://developer.mozilla.org/fr/docs/Glossaire/JavaScript}{Javascript},
    c'est l'un des défauts du langage.
  \item
    Identifier et analyser de même \textbf{le gestionnaire d'événement}
    \passthrough{\lstinline!click!} qui est attaché au second élément
    \passthrough{\lstinline!<button>!}.
  \end{itemize}
\end{enumerate}

\end{exemple}

\begin{exemple}{}

Ouvrir un navigateur Web et accéder à la page
\url{https://replit.com/@fredericjunier/1NSI-Javascript-Ex1-Eleve\#index.html}.

On arrive dans un environnement de développement Web en
\href{https://developer.mozilla.org/fr/docs/Glossaire/HTML}{HTML}/\href{https://developer.mozilla.org/fr/docs/Glossaire/CSS}{CSS}/\href{https://developer.mozilla.org/fr/docs/Glossaire/JavaScript}{Javascript}
sur la plateforme \url{https://repl.it}.

\includegraphics{images/exo4-fig1.png}\\

\begin{enumerate}
\def\labelenumi{\arabic{enumi}.}
\item
  L'interface se divise en trois zones :

  \begin{itemize}
  \item
    À droite la fenêtre d'affichage du site et en-dessous une console
    \href{https://developer.mozilla.org/fr/docs/Glossaire/JavaScript}{Javascript}.
  \item
    Au centre l'éditeur avec le fichier courant.
  \item
    À gauche l'explorateur de fichiers avec une séparation des trois
    composants logiciels nécessaires à l'affichage dans la fenêtre en
    haut à droite :

    \begin{itemize}
    \tightlist
    \item
      la structure du document en
      \href{https://developer.mozilla.org/fr/docs/Glossaire/HTML}{HTML}
      dans \passthrough{\lstinline!index.html!}
    \item
      le code
      \href{https://developer.mozilla.org/fr/docs/Glossaire/JavaScript}{Javascript}
      dans \passthrough{\lstinline!script.js!}
    \item
      la feuille de style
      \href{https://developer.mozilla.org/fr/docs/Glossaire/CSS}{CSS}
      dans \passthrough{\lstinline!style.css!}
    \end{itemize}

    Par rapport à l'exemple précédent, le code
    \href{https://developer.mozilla.org/fr/docs/Glossaire/JavaScript}{Javascript}
    est donc clairement séparé du code
    \href{https://developer.mozilla.org/fr/docs/Glossaire/HTML}{HTML}.
    Le navigateur du client reçoit d'abord le code
    \href{https://developer.mozilla.org/fr/docs/Glossaire/HTML}{HTML}.
    Pour charger le
    \href{https://developer.mozilla.org/fr/docs/Glossaire/JavaScript}{Javascript},
    on procède comme pour une feuille de style
    \href{https://developer.mozilla.org/fr/docs/Glossaire/CSS}{CSS} en
    plaçant une balise de lien dans l'en-tête du fichier
    \href{https://developer.mozilla.org/fr/docs/Glossaire/HTML}{HTML} :

\begin{lstlisting}[language=HTML]
<script src="script.js" defer>
</script>
\end{lstlisting}

    L'attribut \passthrough{\lstinline!src!} donne le chemin vers le
    fichier contenant le code
    \href{https://developer.mozilla.org/fr/docs/Glossaire/JavaScript}{Javascript}
    et l'attribut \passthrough{\lstinline!defer!} précise que le
    chargement du
    \href{https://developer.mozilla.org/fr/docs/Glossaire/JavaScript}{Javascript}
    doit se faire après que la page
    \href{https://developer.mozilla.org/fr/docs/Glossaire/HTML}{HTML}
    soit totalement chargée. En effet,
    \href{https://developer.mozilla.org/fr/docs/Glossaire/JavaScript}{Javascript}
    peut modifier la page à travers l'interface
    \href{https://developer.mozilla.org/fr/docs/Glossaire/DOM}{DOM},
    donc il faut ordonnancer le chargement des différentes ressources.
    Des anciennes pratiques recommandaient ainsi de placer les balises
    \passthrough{\lstinline!<script>!} à la fin du code
    \href{https://developer.mozilla.org/fr/docs/Glossaire/HTML}{HTML}.
    Des attributs comme \passthrough{\lstinline!defer!} permettent de
    gérer plus finement les priorités, surtout lorsque plusieurs scripts
    \href{https://developer.mozilla.org/fr/docs/Glossaire/JavaScript}{Javacript}
    sont appelés.
  \end{itemize}
\item
  On donne ci-dessous le code du fichier
  \passthrough{\lstinline!script.js!} :
\end{enumerate}

\begin{lstlisting}
let compteur = 0;

function suivant(){
    compteur = compteur + 1;
    let v = document.getElementById("valeur");
    if (compteur <= 1) {
      v.innerHTML = compteur + " clic";
    }
    else {
      v.innerHTML = compteur + " clics";
    }          
}

function reinitialiser(){
    compteur = 0;
    let v = document.getElementById("valeur");
    v.innerHTML = compteur + " clic";
}


let bouton_clic = document.getElementById("bouton_clic");
bouton_clic.addEventListener("click", suivant);
let bouton_reset = document.getElementById("bouton_reset");
bouton_reset.addEventListener("click", reinitialiser);
\end{lstlisting}

\begin{itemize}
\item
  Les définitions de la variable \passthrough{\lstinline!compteur!} et
  des gestionnaires d'événements \passthrough{\lstinline!suivant!} et
  \passthrough{\lstinline!reinitialiser!} sont identiques à celles du
  code de l'exemple 1. Ce qui change est la façon dont les gestionnaires
  sont attachés aux boutons. Prenons le bouton de \emph{Clic !}, dans
  l'exemple 1 la liaison se faisait dans la balise
  \href{https://developer.mozilla.org/fr/docs/Glossaire/HTML}{HTML} avec
  l'ajout de l'attribut \passthrough{\lstinline!onclick = "suivant()"!}
  alors qu'ici la liaison est déportée dans le fichier
  \passthrough{\lstinline!script.js!} avec :

\begin{lstlisting}
//récupération de l'élément <button> ciblé par son identifiant et la méthode DOM
let bouton_clic = document.getElementById("bouton_clic");
//ajout sur cet élément du  gestionnaire d'événement suivant pour l'événement click
bouton_clic.addEventListener("click", suivant);
\end{lstlisting}
\item
  L'architecture présentée dans l'exemple 2 est préférable car elle
  respecte un principe de \textbf{séparation des composants logiciels}
  selon leur fonctionnalité et facilite donc la \textbf{lisibilité}, la
  \textbf{factorisation} et la \textbf{maintenance} du code.
\item
  Il faut noter deux autres différences entre les méthodes des exemples
  1 et 2 :

  \begin{itemize}
  \tightlist
  \item
    les parenthèses après \passthrough{\lstinline!suivant!} (qui est une
    fonction) sont présentes dans
    \passthrough{\lstinline!<button onclick = "suivant()">!} mais
    absentes dans
    \passthrough{\lstinline!bouton\_clic.addEventListener("click", suivant);!}
  \item
    la liaison avec l'élément
    \passthrough{\lstinline'<button id="bouton\_clic"> Clic !</button>'}
    se fait dans \passthrough{\lstinline!script.js!} grâce à
    l'identifiant \passthrough{\lstinline!id="bouton\_clic"!} qui
    n'était pas nécessaire dans l'exemple 1 lorsque le gestionnaire
    était attaché directement à la balise.
  \end{itemize}
\end{itemize}

\end{exemple}

\begin{cours}{}

\begin{itemize}
\item
  \href{https://developer.mozilla.org/fr/docs/Glossaire/JavaScript}{Javascript}
  permet grâce à l'interface
  \href{https://developer.mozilla.org/fr/docs/Glossaire/DOM}{DOM}
  d'attacher une fonction \textbf{gestionnaire d'événement} à un élément
  d'un document
  \href{https://developer.mozilla.org/fr/docs/Glossaire/HTML}{HTML}.

  Ce gestionnaire est lié à un \textbf{événement}, qui est une action de
  l'utilisateur (déplacement de la souris, pression sur une touche du
  clavier) ou une modification d'un autre élément du document.

  Lorsque l'événement ciblé atteint l'élément surveillé par le
  gestionnaire, celui-ci le capture et délenche une action.
\item
  Pour attacher un \textbf{gestionnaire d'événement} nommé
  \passthrough{\lstinline!gestionnaire!}, à un élément, nommé
  \passthrough{\lstinline!element!}, et le lier à un événement, par
  exemple \passthrough{\lstinline!click!},on peut procéder de deux
  façons :

  \begin{itemize}
  \tightlist
  \item
    \emph{Méthode 1 (voir exemple 1), directement dans le code
    \href{https://developer.mozilla.org/fr/docs/Glossaire/HTML}{HTML} :}
  \end{itemize}

\begin{lstlisting}[language=HTML]
<element id="id_element" onclick="gestionnaire()"> ... </element>
\end{lstlisting}

  \begin{itemize}
  \tightlist
  \item
    \emph{Méthode 2 (voir exemple 2), dans un fichier externe
    \passthrough{\lstinline!script.js!} :}
  \end{itemize}

\begin{lstlisting}
//la balise element doit être repérée par un id 
let v = document.getElementById("id_element");
v.addEventListener("click", gestionnaire);
\end{lstlisting}
\item
  Il existe plusieurs méthodes pour cibler un élément
  \href{https://developer.mozilla.org/fr/docs/Glossaire/HTML}{HTML},
  nous utiliserons principalement le ciblage par identifiant unique
  \passthrough{\lstinline!id!} :
\end{itemize}

\begin{longtable}[]{@{}cc@{}}
\caption{Différentes façons de cibler un élément
\href{https://developer.mozilla.org/fr/docs/Glossaire/HTML}{HTML}}\tabularnewline
\toprule
\begin{minipage}[b]{0.47\columnwidth}\centering
Type de cible\strut
\end{minipage} & \begin{minipage}[b]{0.47\columnwidth}\centering
Syntaxe\strut
\end{minipage}\tabularnewline
\midrule
\endfirsthead
\toprule
\begin{minipage}[b]{0.47\columnwidth}\centering
Type de cible\strut
\end{minipage} & \begin{minipage}[b]{0.47\columnwidth}\centering
Syntaxe\strut
\end{minipage}\tabularnewline
\midrule
\endhead
\begin{minipage}[t]{0.47\columnwidth}\centering
Identifiant unique \passthrough{\lstinline!id!}\strut
\end{minipage} & \begin{minipage}[t]{0.47\columnwidth}\centering
\passthrough{\lstinline!document.getElementById('id') //retourne un élément!}\strut
\end{minipage}\tabularnewline
\begin{minipage}[t]{0.47\columnwidth}\centering
Classe
\href{https://developer.mozilla.org/fr/docs/Glossaire/CSS}{CSS}\strut
\end{minipage} & \begin{minipage}[t]{0.47\columnwidth}\centering
\passthrough{\lstinline!document.getElementByClass('classname') //retourne une liste d'éléments!}\strut
\end{minipage}\tabularnewline
\begin{minipage}[t]{0.47\columnwidth}\centering
Sélecteur
\href{https://developer.mozilla.org/fr/docs/Glossaire/CSS}{CSS} (par
exemple titre \passthrough{\lstinline!h1!})\strut
\end{minipage} & \begin{minipage}[t]{0.47\columnwidth}\centering
\passthrough{\lstinline!document.querySelector('h1') //retourne une liste d'éléments!}\strut
\end{minipage}\tabularnewline
\bottomrule
\end{longtable}

\begin{itemize}
\tightlist
\item
  Dans le tableau ci-dessous on présente une sélection d'
  \textbf{événements} avec sa description et les deux façons de le lier
  à un gestionnaire \passthrough{\lstinline!gest!} : directement dans
  l'élément
  \href{https://developer.mozilla.org/fr/docs/Glossaire/HTML}{HTML}
  \passthrough{\lstinline!tag!} ou dans un fichier externe où la
  variable \passthrough{\lstinline!v!} désigne cet élément.
\end{itemize}

\begin{longtable}[]{@{}cccc@{}}
\caption{Événements et gestionnaires associés}\tabularnewline
\toprule
\begin{minipage}[b]{0.22\columnwidth}\centering
Événement\strut
\end{minipage} & \begin{minipage}[b]{0.22\columnwidth}\centering
Signification\strut
\end{minipage} & \begin{minipage}[b]{0.22\columnwidth}\centering
Gestionnaire interne\strut
\end{minipage} & \begin{minipage}[b]{0.22\columnwidth}\centering
Gestionnaire externe\strut
\end{minipage}\tabularnewline
\midrule
\endfirsthead
\toprule
\begin{minipage}[b]{0.22\columnwidth}\centering
Événement\strut
\end{minipage} & \begin{minipage}[b]{0.22\columnwidth}\centering
Signification\strut
\end{minipage} & \begin{minipage}[b]{0.22\columnwidth}\centering
Gestionnaire interne\strut
\end{minipage} & \begin{minipage}[b]{0.22\columnwidth}\centering
Gestionnaire externe\strut
\end{minipage}\tabularnewline
\midrule
\endhead
\begin{minipage}[t]{0.22\columnwidth}\centering
\passthrough{\lstinline!click!}\strut
\end{minipage} & \begin{minipage}[t]{0.22\columnwidth}\centering
dispositif de pointage pressé sur l'élément\strut
\end{minipage} & \begin{minipage}[t]{0.22\columnwidth}\centering
\passthrough{\lstinline!<tag onclick="gest()">!}\strut
\end{minipage} & \begin{minipage}[t]{0.22\columnwidth}\centering
\passthrough{\lstinline!v.addEventListener('click', gest)!}\strut
\end{minipage}\tabularnewline
\begin{minipage}[t]{0.22\columnwidth}\centering
\passthrough{\lstinline!mouseover!}\strut
\end{minipage} & \begin{minipage}[t]{0.22\columnwidth}\centering
dispositif de pointage déplacé sur l'élément\strut
\end{minipage} & \begin{minipage}[t]{0.22\columnwidth}\centering
\passthrough{\lstinline!<tag onmouseover="gest()">!}\strut
\end{minipage} & \begin{minipage}[t]{0.22\columnwidth}\centering
\passthrough{\lstinline!v.addEventListener('mouseover', gest)!}\strut
\end{minipage}\tabularnewline
\begin{minipage}[t]{0.22\columnwidth}\centering
\passthrough{\lstinline!mouseout!}\strut
\end{minipage} & \begin{minipage}[t]{0.22\columnwidth}\centering
dispositif de pointage déplacé hors de l'élément\strut
\end{minipage} & \begin{minipage}[t]{0.22\columnwidth}\centering
\passthrough{\lstinline!<tag      onmouseout="gest()">!}\strut
\end{minipage} & \begin{minipage}[t]{0.22\columnwidth}\centering
\passthrough{\lstinline!v.addEventListener('mouseout', gest)!}\strut
\end{minipage}\tabularnewline
\begin{minipage}[t]{0.22\columnwidth}\centering
\passthrough{\lstinline!keydown!}\strut
\end{minipage} & \begin{minipage}[t]{0.22\columnwidth}\centering
une touche du clavier est pressée\strut
\end{minipage} & \begin{minipage}[t]{0.22\columnwidth}\centering
\passthrough{\lstinline!<tag   onkeydown="gest()">!}\strut
\end{minipage} & \begin{minipage}[t]{0.22\columnwidth}\centering
\passthrough{\lstinline!v.addEventListener('keydown', gest)!}\strut
\end{minipage}\tabularnewline
\begin{minipage}[t]{0.22\columnwidth}\centering
\passthrough{\lstinline!keyup!}\strut
\end{minipage} & \begin{minipage}[t]{0.22\columnwidth}\centering
une touche du clavier est pressée\strut
\end{minipage} & \begin{minipage}[t]{0.22\columnwidth}\centering
\passthrough{\lstinline!<tag   onkeyup="gest()">!}\strut
\end{minipage} & \begin{minipage}[t]{0.22\columnwidth}\centering
\passthrough{\lstinline!v.addEventListener('keyup', gest)!}\strut
\end{minipage}\tabularnewline
\begin{minipage}[t]{0.22\columnwidth}\centering
\passthrough{\lstinline!input!}\strut
\end{minipage} & \begin{minipage}[t]{0.22\columnwidth}\centering
à chaque changement de valeur réalisé par l'utilisateur dans
\passthrough{\lstinline!<input>!} ou
\passthrough{\lstinline!<select>!}\strut
\end{minipage} & \begin{minipage}[t]{0.22\columnwidth}\centering
\passthrough{\lstinline!<tag      oninput="gest()">!}\strut
\end{minipage} & \begin{minipage}[t]{0.22\columnwidth}\centering
\passthrough{\lstinline!v.addEventListener('input', gest)!}\strut
\end{minipage}\tabularnewline
\bottomrule
\end{longtable}

\end{cours}

\hypertarget{travaux-pratiques}{%
\section{Travaux pratiques}\label{travaux-pratiques}}

\begin{exercice}{}

Ouvrir un navigateur Web et accéder à la page
\url{https://repl.it/@fredericjunier/1NSI-Javascript-Ex3-Eleve}.

On arrive dans un environnement de développement Web en
\href{https://developer.mozilla.org/fr/docs/Glossaire/HTML}{HTML}/\href{https://developer.mozilla.org/fr/docs/Glossaire/CSS}{CSS}/\href{https://developer.mozilla.org/fr/docs/Glossaire/JavaScript}{Javascript}
sur la plateforme \url{https://repl.it}.

\includegraphics{images/exo-replit.png}\\

La page d'accueil contient une liste de liens vers des pages dynamiques
qui contiennent ou sont liés à des codes
\href{https://developer.mozilla.org/fr/docs/Glossaire/JavaScript}{Javascript}
à compléter. À une exception près, ces activités présentées ont déjà été
implémentés à l'aide de formulaires et de scritps côté serveur dans le
chapitre précédent. Ici l'interactivité sera assurée côté client par
\href{https://developer.mozilla.org/fr/docs/Glossaire/JavaScript}{Javascript}.

\begin{enumerate}
\def\labelenumi{\arabic{enumi}.}
\item
  Depuis la page d'accueil, suivre le lien vers l'activité
  \textbf{Calcul d'âge} de la page \passthrough{\lstinline!age.html!}.

  \begin{itemize}
  \item
    Quel élément
    \href{https://developer.mozilla.org/fr/docs/Glossaire/HTML}{HTM} a
    reçu un gestionnaire d'événement ? Quel événement est surveillé ?
  \item
    Compléter le code
    \href{https://developer.mozilla.org/fr/docs/Glossaire/JavaScript}{Javascript}
    dans le fichier \passthrough{\lstinline!code-age.js!} afin que la
    page puisse calculer l'âge de l'utilisateur à partir de sa date de
    naissance?
  \end{itemize}
\item
  Revenir sur la page d'accueil et suivre le lien vers l'activité
  \textbf{Jouer avec les événements I}. Compléter le code
  \href{https://developer.mozilla.org/fr/docs/Glossaire/JavaScript}{Javascript}
  qui se trouve dans l'en-tête du fichier
  \passthrough{\lstinline!evenements.html!} pour :

  \begin{itemize}
  \item
    attacher un gestionnaire d'événement \passthrough{\lstinline!click!}
    à l'élément d'identifiant \passthrough{\lstinline!titre!}
  \item
    un clic sur cet élément doit déclencher l'apparition d'une fenêtre
    pop-up avec le message ``Attention, je peux griffer !''
  \end{itemize}
\item
  Revenir sur la page d'accueil et suivre le lien vers l'activité
  \textbf{Jouer avec les événements II}. Compléter le code
  \href{https://developer.mozilla.org/fr/docs/Glossaire/JavaScript}{Javascript}
  qui se trouve dans l'en-tête du fichier
  \passthrough{\lstinline!evenements2html!} pour :

  \begin{itemize}
  \item
    attacher un gestionnaire d'événement \passthrough{\lstinline!click!}
    à l'élément \passthrough{\lstinline!img!} d'identifiant
    \passthrough{\lstinline!chat!}
  \item
    un clic sur cet élément doit déclencher le changement de sa
    propriété \passthrough{\lstinline!src!}, qui doit prendre pour
    valeur le chemin ``images/chat-bonjour.png'', ainsi un clic sur
    l'image devra provoquer son changement.
  \end{itemize}
\item
  Revenir sur la page d'accueil et suivre le lien vers l'activité
  \textbf{Colorpicker}. La page \passthrough{\lstinline!rgb.html!}
  propose d'afficher une couleur dans un carré à partir de ses
  composantes (R,G,B). Le code
  \href{https://developer.mozilla.org/fr/docs/Glossaire/JavaScript}{Javascript}
  se trouve dans le fichier \passthrough{\lstinline!code-rgb.js!}.

  \begin{itemize}
  \item
    Quel élément
    \href{https://developer.mozilla.org/fr/docs/Glossaire/HTML}{HTM} a
    reçu un gestionnaire d'événement ? Quel événement est surveillé ?
  \item
    Compléter le code
    \href{https://developer.mozilla.org/fr/docs/Glossaire/JavaScript}{Javascript}
    qui se trouve dans le fichier \passthrough{\lstinline!code-rgb.js!}
    pour rendre la page \passthrough{\lstinline!rgb.html!}
    fonctionnelle.
  \end{itemize}
\item
  Revenir sur la page d'accueil et suivre le lien vers l'activité
  \textbf{Conversion d'unité de température}. La page
  \passthrough{\lstinline!temperature.html!} propose de convertir une
  mesure de température de Celsius en Fahrenheit ou réciproquement. Le
  code
  \href{https://developer.mozilla.org/fr/docs/Glossaire/JavaScript}{Javascript}
  se trouve dans le fichier
  \passthrough{\lstinline!code-temperature.js!}.

  \begin{itemize}
  \item
    Quel élément
    \href{https://developer.mozilla.org/fr/docs/Glossaire/HTML}{HTM} a
    reçu un gestionnaire d'événement ? Quel événement est surveillé ?
  \item
    Compléter le code
    \href{https://developer.mozilla.org/fr/docs/Glossaire/JavaScript}{Javascript}
    qui se trouve dans le fichier
    \passthrough{\lstinline!code-temperature.js!} pour rendre la page
    \passthrough{\lstinline!temperature.html!} fonctionnelle.
  \end{itemize}
\item
  Revenir sur la page d'accueil et suivre le lien vers l'activité
  \textbf{Calcul de la force d'un mot de passe}. La page
  \passthrough{\lstinline!force-mot-passe.html!} propose de tester la
  force d'un mot de passe en fonction de son nombre de caractères et de
  la taille de l'alphabet où sont choisis les caractères.\\
  Le code
  \href{https://developer.mozilla.org/fr/docs/Glossaire/JavaScript}{Javascript}
  se trouve dans le fichier
  \passthrough{\lstinline!code-force-mot-passe.js!}.\\
  Tester d'abord le formulaire de l'ANSSI qui a inspiré cette activité :
  \url{https://www.ssi.gouv.fr/administration/precautions-elementaires/calculer-la-force-dun-mot-de-passe/}.

  \begin{itemize}
  \item
    Quel élément
    \href{https://developer.mozilla.org/fr/docs/Glossaire/HTML}{HTM} a
    reçu un gestionnaire d'événement ? Quel événement est surveillé ?
  \item
    Compléter le code
    \href{https://developer.mozilla.org/fr/docs/Glossaire/JavaScript}{Javascript}
    qui se trouve dans le fichier
    \passthrough{\lstinline!code-force-mot-passe.js!} pour rendre la
    page \passthrough{\lstinline!force-mot-passe.html!} fonctionnelle.
  \end{itemize}
\end{enumerate}

\includegraphics[width=0.6\textwidth,height=\textheight]{images/force-mot-passe.png}\\

\end{exercice}

\hypertarget{qcm-e3c2}{%
\section{QCM E3C2}\label{qcm-e3c2}}

\begin{exercice}{}

\emph{QCM} de type E3C2.

\begin{enumerate}
\def\labelenumi{\arabic{enumi}.}
\item
  On souhaite qu'un menu apparaisse à chaque fois que l'utilisateur
  passe sa souris sur l'image de bannière du site. L'attribut de la
  balise \passthrough{\lstinline!img!} dans lequel on doit mettre un
  code Javascript à cet effet est~:

  \begin{itemize}
  \tightlist
  \item
    Réponse A : \passthrough{\lstinline!onclick!}
  \item
    Réponse B : \passthrough{\lstinline!src!}
  \item
    Réponse C : \passthrough{\lstinline!alt!}
  \item
    Réponse D : \passthrough{\lstinline!onmouseover!}
  \end{itemize}
\item
  Lors de la consultation d'une page HTML contenant un bouton auquel est
  associée la fonction suivante, que se passe-t-il quand on clique sur
  ce bouton ?

\begin{lstlisting}
function action(event) {
  this.style.color = "blue"
  }
\end{lstlisting}

  \begin{itemize}
  \tightlist
  \item
    Réponse A : le texte de la page passe en bleu
  \item
    Réponse B : le texte du bouton passe en bleu
  \item
    Réponse C : le texte du bouton est changé et affiche maintenant le
    mot ``bleu''
  \item
    Réponse D : le pointeur de la souris devient bleu quand il arrive
    sur le bouton
  \end{itemize}
\item
  Parmi les propriétés suivantes d'une balise
  \passthrough{\lstinline!<button>!} dans une page HTML, laquelle doit
  être rédigée en langage JavaScript ?

  \begin{itemize}
  \tightlist
  \item
    Réponse A : la propriété \passthrough{\lstinline!name!}
  \item
    Réponse B : la propriété \passthrough{\lstinline!type!}
  \item
    Réponse C : la propriété \passthrough{\lstinline!onclick!}
  \item
    Réponse D : la propriété \passthrough{\lstinline!id!}
  \end{itemize}
\item
  Dans une page HTML, lequel de ces codes permet la présence d'un bouton
  qui appelle la fonction javascript
  \passthrough{\lstinline!afficher\_reponse()!} lorsque l'utilisateur
  clique dessus ?

  \begin{itemize}
  \tightlist
  \item
    Réponse A :
    \passthrough{\lstinline!<a href="afficher\_reponse()">Cliquez ici</a>!}
  \item
    Réponse B :
    \passthrough{\lstinline!<button if\_clicked="afficher\_reponse()">Cliquez ici</button>!}
  \item
    Réponse C :
    \passthrough{\lstinline!<button value="Cliquez ici"><a>afficher\_reponse()</a></button>!}
  \item
    Réponse D :
    \passthrough{\lstinline!<button onclick="afficher\_reponse()">Cliquez ici</button>!}
  \end{itemize}
\item
  Quel est le nom de l'événement généré lorsque l'utilisateur clique sur
  un bouton de type button dans une page HTML ?

  \begin{itemize}
  \tightlist
  \item
    Réponse A : \passthrough{\lstinline!action!}
  \item
    Réponse B : \passthrough{\lstinline!mouse!}
  \item
    Réponse C : \passthrough{\lstinline!submit!}
  \item
    Réponse D : \passthrough{\lstinline!click!}
  \end{itemize}
\item
  Un navigateur affiche la page HTML suivante :

\begin{lstlisting}[language=HTML]
<html lang="fr">
<head>
<meta charset="utf-8">
<link rel="stylesheet" href="style.css">
<title>Un bouton</title>
</head>
<body>
<button onclick="maFonction()">Cliquer ici</button>
</body>
<script src="script.js"></script>
</html>
\end{lstlisting}

  Lorsque l'on clique sur le bouton, l'action déclenchée par
  \passthrough{\lstinline!maFonction()!} est définie :

  \begin{itemize}
  \tightlist
  \item
    Réponse A : dans le fichier HTML seul
  \item
    Réponse B : dans le fichier \passthrough{\lstinline!style.css!}
  \item
    Réponse C : dans une bibliothèque prédéfinie du navigateur
  \item
    Réponse D : dans le fichier \passthrough{\lstinline!script.js!}
  \end{itemize}
\item
  Voici un extrait d'une page HTML :

\begin{lstlisting}[language=HTML]
<script>
function sommeNombres(formulaire) {
var somme = formulaire.n1.value + formulaire.n2.value;
console.log(somme);
}
</script>

<form>
Nombre 1 : <input name="n1" value="30"> <br>
Nombre 2 : <input name="n2" value="10"> <br>
<input type="button" value="Somme"
onclick="sommeNombres(this.form)">
</form>
\end{lstlisting}

  Quand l'utilisateur clique sur le bouton Somme, le calcul de la
  fonction \passthrough{\lstinline!sommeNombre()!} se fait~:

  \begin{itemize}
  \tightlist
  \item
    Réponse A : uniquement dans le navigateur
  \item
    Réponse B : uniquement sur le serveur qui héberge la page
  \item
    Réponse C : à la fois dans le navigateur et sur le serveur
  \item
    Réponse D : si le calcul est complexe, le navigateur demande au
    serveur de faire le calcul
  \end{itemize}
\item
  On considère cet extrait de fichier HTML représentant les onglets
  d'une barre de navigation :

\begin{lstlisting}
function BoutonGris() {

var btn = document.createElement(\"BUTTON\");

btn.innerHTML = \"Annulation\";

document.getElementById(\"DIV\").appendChild(btn);

}
\end{lstlisting}

  \begin{itemize}
  \tightlist
  \item
    Réponse A : elle remplace un élément DIV par un bouton
  \item
    Réponse B : elle annule l'élément BUTTON
  \item
    Réponse C : elle crée un bouton comportant le texte ``Annulation''
  \item
    Réponse D : elle recherche le bouton ``BUTTON'' et crée une copie
    appelée ``btn''
  \end{itemize}
\end{enumerate}

\end{exercice}

\begin{figure}
\centering
\includegraphics[width=0.25\textwidth,height=\textheight]{images/types.png}
\caption{``\href{https://xkcd.com/1537/}{XKCD 1537 : Types}''}
\end{figure}

\end{document}
